% !TEX root = ../main.tex
%% Devoir surveillé n° 2 — chapitres 4 à 7 (valeurs vérifiées : verifications/devoirs.py)

\devoirDS{2}{Chapitres 4 à 7}{{L'eau oxygénée à 10 volumes : titrage colorimétrique}/8,{Trois molécules : schémas de Lewis, géométrie, polarité}/6,{Le limonène de l'orange : extraction par un solvant}/6}

\begin{exoBac}{8}{L'eau oxygénée à 10 volumes}
L'eau oxygénée est une solution aqueuse de peroxyde d'hydrogène $\ce{H2O2}$, utilisée comme antiseptique. Une solution « à 10
volumes » peut libérer 10 litres de dioxygène par litre de solution, mesurés à \qty{0}{\degreeCelsius} sous \qty{1,013}{bar}, selon la réaction
$\ce{2 H2O2 -> 2 H2O + O2}$. On contrôle cette indication par un titrage.
\begin{docu}{1}{Le protocole}
On dilue 10 fois la solution commerciale : on obtient la solution $\mathrm{S_1}$. On titre $V_1=\qty{10,0}{mL}$ de $\mathrm{S_1}$, acidifiée, par une
solution de permanganate de potassium, $\ce{K+}+\ce{MnO4-}$, de concentration $c_2=\qty{2,00e-2}{mol.L^{-1}}$. L'équivalence est obtenue
pour $V_\mathrm{E}=\qty{17,8}{mL}$.
\end{docu}
\begin{donnees}
Couples oxydant-réducteur : $\ce{MnO4-}/\ce{Mn^{2+}}$ et $\ce{O2}/\ce{H2O2}$. En solution, les ions $\ce{MnO4-}$ sont violets ; les ions
$\ce{Mn^{2+}}$ et le peroxyde d'hydrogène sont incolores. Volume molaire des gaz à \qty{0}{\degreeCelsius} sous \qty{1,013}{bar} : \qty{22,4}{L.mol^{-1}}.
Burette graduée de \qty{25}{mL}.
\end{donnees}
\begin{questions}
\item Écrire les demi-équations électroniques des deux couples, puis montrer que l'équation de la réaction support du titrage
s'écrit : $\ce{2 MnO4- + 5 H2O2 + 6 H+ -> 2 Mn^{2+} + 5 O2 + 8 H2O}$.
\item Faire un schéma légendé du montage de titrage.
\item Comment repère-t-on l'équivalence ? Pourquoi n'est-il pas nécessaire d'ajouter un indicateur coloré ?
\item Définir l'équivalence et établir la relation entre la quantité de peroxyde d'hydrogène titrée et la quantité d'ions
permanganate versée à l'équivalence.
\item Calculer la concentration en peroxyde d'hydrogène de $\mathrm{S_1}$, puis celle de la solution commerciale.
\item \initiative L'indication « 10 volumes » est-elle vérifiée ?
\item Pourquoi a-t-on dilué la solution commerciale avant de la titrer ?
\end{questions}
\end{exoBac}

\begin{exoBac}{6}{Trois molécules}
On étudie trois molécules : le méthanal $\ce{CH2O}$, utilisé dans les colles et les résines ; le cyanure d'hydrogène $\ce{HCN}$, gaz très
toxique ; le tétrachlorométhane $\ce{CCl4}$, ancien solvant aujourd'hui interdit. Dans ces molécules, l'atome de carbone est central.
\begin{donnees}
Numéros atomiques : H 1 ; C 6 ; N 7 ; O 8 ; Cl 17. Électronégativités (échelle de Pauling) : H 2,20 ; C 2,55 ; N 3,04 ; Cl 3,16 ;
O 3,44. On considère qu'une liaison est polarisée si l'écart d'électronégativité entre ses atomes dépasse 0,4.
\end{donnees}
\begin{questions}
\item Écrire la configuration électronique des atomes de carbone, d'azote, d'oxygène et de chlore. En déduire leur nombre
d'électrons de valence.
\item Établir le schéma de Lewis de chaque molécule.
\item Prévoir la géométrie de chaque molécule.
\item Repérer les liaisons polarisées et indiquer les charges partielles.
\item Quelles molécules sont polaires ? Justifier en particulier le cas du tétrachlorométhane.
\item Le méthanal est très soluble dans l'eau, contrairement au tétrachlorométhane. Interpréter.
\end{questions}
\end{exoBac}

\begin{exoBac}{6}{Le limonène de l'orange}
Le limonène, $\ce{C10H16}$, est la molécule responsable de l'odeur de l'orange ; c'est une molécule apolaire. Par hydrodistillation
d'écorces d'orange, on obtient un distillat formé d'eau et de limonène, que l'on veut extraire par un solvant.
\begin{docu}{1}{Quelques solvants}
{\small\begin{tabularx}{\linewidth}{@{}lYYYY@{}}\toprule
 & eau & éthanol & cyclohexane & dichlorométhane \\\midrule
Miscibilité avec l'eau & -- & totale & nulle & nulle \\
Densité & 1,00 & 0,79 & 0,78 & 1,33 \\
Solubilité du limonène & très faible & grande & grande & grande \\
Sécurité & -- & inflammable & inflammable, nocif pour l'environnement & suspecté d'être cancérogène \\\bottomrule
\end{tabularx}}
\end{docu}
\begin{donnees}
Températures d'ébullition : cyclohexane \qty{81}{\degreeCelsius} ; limonène \qty{176}{\degreeCelsius}.
\end{donnees}
\begin{questions}
\item Pourquoi le limonène est-il très peu soluble dans l'eau ?
\item Quel solvant ne peut pas convenir pour l'extraction ? Justifier.
\item Choisir, en justifiant, entre le cyclohexane et le dichlorométhane.
\item Décrire l'extraction à l'aide d'une ampoule à décanter. Schématiser l'ampoule après décantation en indiquant la nature et la
position de chaque phase.
\item Comment récupère-t-on la phase qui contient le limonène ? Comment obtenir ensuite le limonène seul ?
\end{questions}
\end{exoBac}
